% !PRINTED 595
\weberpage{595}
%<<CONT
woraus sich ergiebt, dass \gesperrt{jede Substitution $\psi$ aus $\Psi$ in einem der Systeme} (Nebengruppen)
\[ X,\ X\psi_0,\ X\psi_0^2,\ \ldots \]
\gesperrt{enthalten ist.}

Wenn $\mathfrak{p}$ vom Grade $f$ (in Bezug auf $R$) ist, so ist
\[ \omega^{P^f} \equiv \omega\ (\mathrm{mod}\ \mathfrak{p}), \]
und $P^f$ ist die niedrigste Potenz von $P$ mit positiven Exponenten, die dieser Congruenz für alle $\omega$ genügt.

Setzt man also $\nu = f$ in (10), so folgt, dass $\psi_0^f$ in $X$ enthalten ist, dass aber keine Potenz von $\psi_0$ mit niedrigerem positiven Exponenten diese Eigenschaft hat. Die Nebengruppen
\[ X,\ X\psi_0,\ X\psi_0^2,\ \ldots,\ X\psi_0^{f-1} \]
sind also alle von einander verschieden, und es ergiebt sich:
\weq{13}{\Psi = X + X\psi_0 + X\psi_0^2 + \cdots + X\psi_0^{f-1}.}

Nach (12) ist aber die Gesammtheit der Substitutionen $\psi_0^\nu X$ mit $X\psi_0^\nu$ identisch (für jedes $\nu$), also
\weq{14}{\psi_0^\nu X = X\psi_0^\nu,}
woraus folgt, dass \gesperrt{$X$ ein Normaltheiler von $\Psi$ ist.}

\gesperrt{Da, wie wir oben gesehen haben, $\Psi$ vom Grade $gf$ ist, so ist $X$ vom Grade $g$.}

In der Gruppe $X$ giebt es nun solche Substitutionen $\chi_1$, darunter immer die identische, für die die Congruenz
\weq{15}{\tau\mkern1mu|\mkern1mu\chi_1 \equiv \tau\ (\mathrm{mod}\ \mathfrak{p}^2)}
erfüllt ist, und diese bilden eine Gruppe $X^{(1)}$, deren Grad mit $g_1$ bezeichnet sein mag.

\gesperrt{Es ist zu beweisen, dass $g_1$ eine Potenz von $p$ ist.} Bezeichnen wir nämlich mit $\pi$ ein durch $\mathfrak{p}$ theilbares Primfunctional, mit $\alpha$ ein ganzes Functional, so folgt aus (15):
\weq{16}{\tau\mkern1mu|\mkern1mu\chi_1 = \tau + \pi^2\alpha,}
und da ebenso nach (15)
\[ \alpha\mkern1mu|\mkern1mu\chi_1 \equiv \alpha\ (\mathrm{mod}\ \mathfrak{p}^2),\quad (\pi\mkern1mu|\mkern1mu\chi_1)^2 \equiv \pi^2\ (\mathrm{mod}\ \mathfrak{p}^3) \]
ist, so ergiebt sich aus (16)
\[ \tau\mkern1mu|\mkern1mu\chi_1^2 \equiv \tau + 2\pi^2\alpha\ (\mathrm{mod}\ \mathfrak{p}^3), \]
und allgemein für jeden positiven Exponenten $h$
\[ \tau\mkern1mu|\mkern1mu\chi_1^h \equiv \tau + h\pi^2\alpha\ (\mathrm{mod}\ \mathfrak{p}^3). \]
